\subsection{Projective three-space concordances and odd local homeomorphisms}

A concordance of real projective three-space lifts to an odd
homeomorphism of an annulus.  Filling its inner boundary by the
identity relates the concordance space to the group of odd local
homeomorphisms at the origin.  The distinction between extending
one local map and extending a whole parameter family is central
to this description.

Let $D=D^4$ be the closed Euclidean unit ball, let $S=\partial D$,
and let $\iota(x)=-x$.  Give
\[
 H=\operatorname{Homeo}_{\iota}(D),\qquad
 A=\operatorname{Homeo}_{\iota}(S)
\]
their compact-open topologies.  A member of $H$ fixes $0$, the
unique fixed point of $\iota$.  Define the simplicial group
$\mathcal G_{\iota}$ of \emph{joint odd germs} as follows.
A $q$-simplex is represented by a continuous, parameter-preserving
homeomorphism between neighborhoods of
$\Delta^q\times\{0\}$ in $\Delta^q\times\mathbb R^4$,
fixing that zero section and commuting with $\iota$ in each fibre.
Two representatives are equal when they agree on a neighborhood
of the entire zero section.  Compactness of $\Delta^q$ gives
common smaller balls for maps and inverses.  Faces and
degeneracies are parameter pullbacks.  Write $\mathcal G_4$
for the same construction with unrestricted local homeomorphisms.

An arrow from a topological group $X$ to a geometric realization
below denotes the derived map represented by the evaluation weak
equivalence $|\operatorname{Sing}X|\to X$ and the indicated realized
simplicial homomorphism.  Its homotopy fibre is computed from this
span.  In particular, the cone-germ map is represented by
\[
 A\ \xleftarrow{\simeq_w}\ |\operatorname{Sing}A|
 \longrightarrow |\operatorname{Sing}H|
 \xrightarrow{|q|}\ |\mathcal G_{\iota}|.
\]

\begin{theorem}\label{thm:projective-concordance-joint-germ}
Restriction to a joint germ defines a Kan fibration
\[
 q:\operatorname{Sing}H\longrightarrow\mathcal G_{\iota}
\]
which is onto vertices.  Its identity fibre is weakly equivalent
to $\operatorname{Sing}P_{\mathrm{TOP}}(\mathbb{RP}^3)$.
Radial coning and the linear inclusion give weak equivalences
\begin{equation}\label{eq:projective-concordance-germ-fibre}
 \begin{split}
 P_{\mathrm{TOP}}(\mathbb{RP}^3)
 &\simeq_w\operatorname{hofib}_1
       \bigl(A\longrightarrow|\mathcal G_{\iota}|\bigr)\\
 &\simeq_w\operatorname{hofib}_1
       \bigl(\mathrm O(4)\longrightarrow|\mathcal G_{\iota}|\bigr)\\
 &\simeq_w\Omega\bigl|\mathcal G_{\iota}/
                                   \operatorname{Sing}\mathrm O(4)\bigr|.
 \end{split}
\end{equation}
All orientation components are included.  Forgetting oddness
induces an isomorphism of groups
\begin{equation}\label{eq:projective-concordance-forgetful-kernel}
 \pi_0P_{\mathrm{TOP}}(\mathbb{RP}^3)
 \cong
 \ker\!\left(\pi_1|\mathcal G_{\iota}|
                 \longrightarrow\pi_1|\mathcal G_4|\right).
\end{equation}
In particular, the full concordance component group is abelian.
\end{theorem}

\begin{proof}
The geometric existence input is the equivariant annulus theorem
of Boyle--Chen--Conway: every free involution on $S^3\times I$
preserving its two boundary components is conjugate to the
antipodal product involution
\cite[Theorem~2.3, p.~6]{BoyleChenConway2025}.
Their proof identifies the quotient as a projective-space
$h$-cobordism and applies the topological product theorem of
Kwasik--Schultz
\cite[Theorem~2.1, p.~530; proof pp.~530--532]{KwasikSchultz1989}.
The ordinary four-dimensional annulus theorem supplies the
underlying annulus, as in
\cite[Theorem~1.4, p.~15]{BoyleChenConway2025}.
For the parameter families we construct a continuous local
extension operator from Edwards--Kirby's deformation theorem.

\paragraph{A normalized extension operator near the identity.}
Choose $0<a<b<c<d<1$.  In the punctured disk let
$C=\{b\leq\|x\|\leq c\}$, and choose a compact annulus
$U$ with $C\subset\operatorname{int}U\subset
U\subset\{a<\|x\|<d\}$.  Its free quotient lies in
\[
 M=(\operatorname{int}D\setminus\{0\})/\iota
           \cong\mathbb{RP}^3\times(0,1).
\]
An odd local homeomorphism $f$, defined on a ball containing
$U$, fixing $0$, and sufficiently close to the identity,
induces an embedding $\bar f:\bar U\to M$ near the inclusion.
We restrict the neighborhood of $f$ so that the relevant
closed source ball has image in $\operatorname{int}D$.

Edwards--Kirby's Theorem~5.1(1) deforms $\bar f$ through
embeddings $\varphi(\bar f,t)$, unchanged near $\partial\bar U$,
to an embedding equal to the identity on $\bar C$.  The
deformation fixes the inclusion and depends continuously on
$\bar f$ \cite[p.~70]{EdwardsKirby1971}.
The images of $\bar U$ are constant during this deformation:
the boundary is fixed, and invariance of domain makes the
membership of every point off that boundary constant along
the isotopy.  Thus the formula
\[
 \bar k_{f,t}=\bar f\,\varphi(\bar f,t)^{-1}
 \quad\hbox{on }\bar f(\bar U),
 \qquad \bar k_{f,t}=1\quad\hbox{elsewhere}
\]
defines an ambient isotopy.  Near the joining boundary the
two formulas agree.  This is the ambient correction in the
proof of \cite[Corollary~1.2, p.~79]{EdwardsKirby1971}.
After shrinking the embedding neighborhood, its support lies
in one compact annulus.  It starts at the identity, its
endpoint agrees with $\bar f$ on $\bar C$, and the entire
isotopy is the identity when $f$ is the identity.

Lift this isotopy through the double cover, starting at the
identity.  The lift $k_{f,t}$ is odd, is identical near $0$
and $\partial D$, and extends there by the identity.
Covering homotopy gives continuous dependence on $f$ and $t$.
Its endpoint $k_f$ agrees with $f$ on $C$: the two lifts
differ there by a constant deck transformation, and closeness
to the identity at one point of the connected collar selects
the identity deck transformation.  Fix $r\in(b,c)$ and set
\begin{equation}\label{eq:projective-normalized-local-extension}
 E(f)(x)=
 \begin{cases}
  f(x),&\|x\|\leq r,\\
  k_f(x),&\|x\|\geq r.
 \end{cases}
\end{equation}
The two maps agree on an annular neighborhood of the joining
sphere.  They map that sphere to the same locally flat sphere,
and both inner regions are the component containing $0$.
Consequently \eqref{eq:projective-normalized-local-extension}
is an odd homeomorphism of $D$.  Pasting gives joint evaluation
continuity, and inversion in the compact-domain homeomorphism
group gives joint inverse continuity.  This defines a single
continuous operator, equal to $f$ near $0$, with $E(1)=1$.

\paragraph{Extension of individual germs.}
Given an odd germ $f$, choose $\epsilon>0$ so that $f$ is
defined near $D_\epsilon$ and
$f(D_\epsilon)\subset\operatorname{int}D$.
The ordinary annulus theorem identifies
$D\setminus\operatorname{int}f(D_\epsilon)$ with $S^3\times I$.
The inherited involution is free and preserves both boundaries.
The equivariant annulus theorem therefore gives an odd product
parametrization.  Precompose its sphere coordinate with the
equivariant boundary homeomorphism needed to agree with $f$
on the inner sphere, using that same correction on every
product level.  Gluing this product to $f$ on $D_\epsilon$
gives an element of $H$ representing the germ.  This proves
surjectivity on vertices.

\paragraph{Relative parameter lifting.}
Let $T=\Delta^j$, let $Z\subset T$ be a horn, and suppose
a joint germ $g$ over $T$ has a prescribed continuous
$H$-lift $F_z$ on $Z$.  The finitely many faces paste to
one family on $Z$.  Choose a retraction $r:T\to Z$ and
a deformation $r_t:T\to T$ from $r$ to the identity,
fixed on $Z$.  Put $\widehat F_p=F_{r(p)}$ and, on a
common sufficiently small ball, put
\[
 \delta_p=\widehat F_p^{-1}g_p,\qquad
 \delta_{p,t}=\delta_{r_t(p)}.
\]
The joint-germ equality on each face, followed by the
finite intersection of the resulting neighborhoods, makes
$\delta_z=1$ on one common ball for all $z\in Z$.
Thus $\delta_{p,0}=1$ for all $p$, and
$\delta_{z,t}=1$ for all $(z,t)\in Z\times I$.

Compactness and joint inverse continuity give a common
ball on which all increments
\[
 \eta_{p,t}=\delta_{p,t}\delta_{p,t_i}^{-1},
                       \qquad t_i\leq t\leq t_{i+1},
\]
are defined.  Use a fixed smaller annular collar there
for a rescaled copy of the operator $E$.
Uniform continuity near the diagonal in the compact
parameter space of $(p,t_i,t)$ supplies a finite time
partition so that each increment belongs to its domain.
Choose a further common ball $V$ whose images under all
$\delta_{p,t}$ lie in the inner ball where $E$ agrees
with its argument.  Starting with $D_{p,0}=1$, define
successively
\[
             D_{p,t}=E(\eta_{p,t})D_{p,t_i}.
\]
The formulas agree at subdivision times because $E(1)=1$.
Induction shows that $D_{p,t}=\delta_{p,t}$ on $V$;
the image condition on $V$ makes each inductive composition
valid on this same ball.  On $Z$ every factor is the
identity.  Hence $\widehat F_pD_{p,1}$ lifts $g$ and
agrees exactly with the prescribed horn family.
This proves the Kan condition.

\paragraph{The identity fibre.}
Let $\mathcal F$ be the strict identity fibre.  Each of
its simplices is an actual $H$-family equal to the
identity on a common ball $D_\epsilon$.  Fix
$0<\epsilon_0<1$, and write $H_{\epsilon_0}$ for the
subgroup fixing $D_{\epsilon_0}$ pointwise.  Taking
the antipodal quotient of the outer annulus identifies
this group homeomorphically with
$P_{\mathrm{TOP}}(\mathbb{RP}^3)$.
Conversely a bottom-fixed projective concordance lifts
uniquely with identity lift on the inner sphere and
extends by identity across the ball.  Covering lifts
and inverse evaluations are continuous for compact
parameter families.  These identifications respect
composition.

The inclusion $\operatorname{Sing}H_{\epsilon_0}\to
\mathcal F$ is a weak equivalence.  A finite simplicial
representative or homotopy in $\mathcal F$ uses one
$\epsilon\leq\epsilon_0$.  Let $a_t$ be the odd radial
homeomorphism of $D$ whose radius function is affine
between the three points
\[
 (0,0),\quad
 (\epsilon,(1-t)\epsilon+t\epsilon_0),\quad(1,1).
\]
It fixes the outer sphere and expands all radii.
Conjugation by $a_t$ moves the given family into
$H_{\epsilon_0}$ at $t=1$.  It preserves that subgroup
throughout because $a_t(D_{\epsilon_0})\supset
D_{\epsilon_0}$.  This proves surjectivity on homotopy
groups.  Applied to a null-homotopy, the same conjugation
gives a null-homotopy in $H_{\epsilon_0}$ after attaching
its boundary conjugation track, which lies in that
subgroup.  This proves injectivity.  The argument also
treats components and keeps the identity basepoint fixed.

\paragraph{Replacing disk maps by orthogonal maps.}
Boundary restriction $H\to A$ has the homomorphic
section $c(a)(ru)=r a(u)$.  Its kernel contracts by
the odd Alexander formula
\[
 \mathcal A_t(h)(x)=
 \begin{cases}
  t h(x/t),&\|x\|\leq t,\quad t>0,\\
  x,&\|x\|\geq t,
 \end{cases}
 \qquad\mathcal A_0(h)=1.
\]
Maps and inverses are continuous in all compact
parameters.  The section gives a product decomposition
of the underlying spaces into this kernel and $A$,
so restriction is a homotopy equivalence.
On the section, $q$ is precisely the cone-germ map.
The Kan fibre sequence now proves the first equivalence
in \eqref{eq:projective-concordance-germ-fibre}.

Bamler--Kleiner prove that
$\operatorname{Isom}(N)\to\operatorname{Diff}^{\infty}(N)$
is a homotopy equivalence for every round spherical
three-manifold, including $N=\mathbb{RP}^3$
\cite[Theorem~1.3, p.~2; proof p.~108]{BamlerKleiner2019Metrics}.
Cerf's comparison of diffeomorphisms with homeomorphisms
applies to compact three-manifolds; its sphere Smale
hypothesis is supplied by that same theorem
\cite[III, \S3.2.1, Theorem~8, p.~365, and Corollary~1,
p.~366; proof pp.~367--370]{Cerf1961}.

We spell out the change of differentiability convention
in this application.  Cerf permits the $C^1$ topology,
whereas the preceding diffeomorphism group uses the
$C^{\infty}$ topology.  For a closed smooth manifold,
the inclusion of these two diffeomorphism groups is
a weak equivalence.  Indeed, embed the manifold in
Euclidean space and fix a smooth retraction of a
tubular neighborhood.  Given a compact parameter
family of $C^1$ diffeomorphisms, choose finitely many
nearby smooth maps by chartwise convolution and a
spatial partition of unity.  A parameter partition
of unity forms finite linear combinations of these
maps.  Composing with the tubular retraction produces
a family continuous in the $C^{\infty}$ topology.
It can be made uniformly $C^1$-close to the original
family.  The diffeomorphisms form an open subset of
the $C^1$ mapping space, so compactness gives a
uniform neighborhood in which both this family and
the retracted straight-line interpolation consist
of diffeomorphisms.

For a sphere or disk parameter pair, first make the
original family constant in the inward normal
parameter near its boundary.  If its boundary family
is already continuous in $C^{\infty}$, use that given
family unchanged on a smaller parameter collar and
interpolate to the finite approximation away from
the collar.  The interpolation weights vanish on
the smaller collar, so this approximation and its
$C^1$ homotopy are relative to the boundary.  The
given boundary family has bounded derivatives of
each order on its compact parameter set; hence the
pasted approximating family is continuous in every
smooth seminorm.  Relative approximation proves
surjectivity and injectivity on all homotopy groups,
including the component assertion.  This establishes
the stated change of convention for $\mathbb{RP}^3$.

Every homeomorphism of $\mathbb{RP}^3$ has two lifts
to $S^3$, both commuting with the deck involution.
The lift space is a double cover of its homeomorphism
group: close maps have a unique close lift after
one sheet is specified.  The same applies to
diffeomorphisms and isometries.  The compatible
diagram of double covers over
\[
 \mathrm{PO}(4)\longrightarrow
 \operatorname{Diff}^{\infty}(\mathbb{RP}^3)
 \longrightarrow\operatorname{Homeo}(\mathbb{RP}^3)
\]
therefore gives a weak equivalence $\mathrm O(4)\to A$.
Its cone-germ map is the literal linear-germ inclusion.
This proves the second equivalence in
\eqref{eq:projective-concordance-germ-fibre}.
Finally, the simplicial subgroup
$\operatorname{Sing}\mathrm O(4)\subset\mathcal G_{\iota}$
has a Kan coset fibration.  The homotopy fibre of its
fibre inclusion is the loop space of that coset space,
which proves the last equivalence.

\paragraph{Components and the forgetful kernel.}
Vertex surjectivity and $H\simeq A\simeq_w\mathrm O(4)$
show that $\pi_0\mathrm O(4)\to\pi_0\mathcal G_{\iota}$
is onto.  The local orientation character separates
the positive and negative linear germs, making it
bijective.  The exact sequence of the simplicial
group fibration gives an isomorphism of groups
\[
 \pi_0P_{\mathrm{TOP}}(\mathbb{RP}^3)
 \cong\operatorname{coker}\bigl(
     \pi_1\mathrm O(4)\longrightarrow
     \pi_1|\mathcal G_{\iota}|\bigr).
\]
The group law on the left is composition, by the
fixed-radius fibre identification.  Fundamental
groups of simplicial groups are abelian, so this
cokernel is an abelian group.

Put $V=\pi_1|\mathcal G_{\iota}|$,
$W=\pi_1|\mathcal G_4|$, and let
\[
 j:\pi_1\mathrm O(4)\longrightarrow V,
 \qquad f:V\longrightarrow W
\]
be the linear inclusion and forgetting map.
The global-to-joint-germ comparison in the proof
of Theorem~\ref{thm:top-sphere-finite-rank-stabilization}
identifies $W$ with $\pi_1\mathrm{TOP}(4)$ while
preserving the literal linear map.  The stabilization
range used in Corollary~\ref{cor:four-sphere-low-concordance-groups}
gives $\pi_1Y_4=\pi_2Y_4=0$ for the smoothing fibre
\eqref{eq:finite-rank-top-smoothing-space}.
The sequence $\mathrm O(4)\to\mathrm{TOP}(4)\to Y_4$
therefore makes $fj$ an isomorphism.  Since
$\pi_1\mathrm O(4)=\mathbb Z/2$, it follows that
\[
 V=j(\mathbb Z/2)\oplus\ker f,
 \qquad
 V/j(\mathbb Z/2)\xrightarrow{\ \cong\ }\ker f,
 \quad [v]\longmapsto v-j\bigl((fj)^{-1}f(v)\bigr).
\]
Combining these isomorphisms proves
\eqref{eq:projective-concordance-forgetful-kernel}.
\end{proof}

Let $K$ denote the kernel in
\eqref{eq:projective-concordance-forgetful-kernel}, and
let $I\subset K$ be the image of the PL concordance
component group.  Transporting the Casson--Sullivan
homomorphism gives an epimorphism
$\chi:K\to\mathbb F_2$ that vanishes on $I$, by
Theorems~\ref{thm:natural-three-concordance-obstruction}
and~\ref{thm:cyclic-three-concordance-obstruction}.
Since $K$ is abelian, this gives the exact sequence
\[
 0\longrightarrow\ker\chi/I\longrightarrow K/I
             \xrightarrow{\ \chi\ }\mathbb F_2\longrightarrow0.
\]
Determining $K$, its PL subgroup $I$, and the remaining
quotient $\ker\chi/I$ is the resulting comparison problem.
